\scanpage{12}
$(\varphi,\omega)$ ($\varphi$ forme bilinéaire sur $E$, $\omega$ base de $\underline{\det} E$) se trouvent multipliés par un $=$ facteur $\lambda$, quand $\omega$ est multiplié par $\lambda$\inserted{ et $\varphi$ \unsure{reste|restant} \unsure{fixée|fixe}}. D'autre part, pour $\omega$ fixé et $\varphi$ variable, le crochet $[\, ,\,]_{\varphi\deleted{\unsure{'}}}$ sur $E'$ dépend linéairement de $\varphi$, tandis que le crochet $[\, ,\,]_{\varphi'}$ sur $E$ (tout comme $f'$ lui-$=$ dont il est déduit linéairement) dépend \uline{quadratiquement} de $\varphi$ ; ainsi, si $\varphi$ est \deleted{multiplié}\inserted{remplacé} par $\mu\varphi$ ($\mu\in\Gamma(S,\underline{\mathrm O}_S)$) $[\, ,\,]_\varphi$ est multiplié par $\mu$, $[\, ,\,]_{\varphi'}$ par $\mu^2$....

Échangeant les rôles de $E,E'$, on voit de $=$ qu'à $f':E'\to E$ est associée une $(f')'=f'':E\to E'$ -- mais en général celle-ci \uline{n'est pas} $f$. De façon précise, si \deleted{$\delta(\varphi')=$}
\[
 \delta(\varphi')=\Delta'\mathbin{.}\omega^{\otimes2}\qquad
 \Delta'\in\Gamma(\unsure{S'|S},\underline{\mathrm O}_{\unsure{S'|S}}) \tag*{(43)}
\]
on aura
\[
 \left\{\begin{aligned}
 f'f''&=\Delta'\,\mathrm{id}_E\\
 f''f'&=\Delta'\,\mathrm{id}_{E'}
 \end{aligned}\right. \tag*{(44)}
\]
ce qui \deleted{\unsure{sugg-}} implique, si $\Delta$ est une unité i.e. $f$ un iso, donc $\bigwedge^2f$ et $f'$ aussi, donc $\Delta'$ une unité [et inversement, car si $f'$ est un iso les formules (42)

% archive sheet number: 11
